\documentclass[11pt,a4paper]{article}
\usepackage[margin=26mm]{geometry}
\usepackage{amsmath,amssymb,amsthm,mathtools}
\usepackage{lmodern,microtype,booktabs,longtable,array}
\usepackage[colorlinks=true,linkcolor=blue,citecolor=blue,urlcolor=blue]{hyperref}
\hypersetup{pdftitle={Optimizing compensated-probe geometry below seven eighths},pdfauthor={RH-Weil research project}}
\newtheorem{theorem}{Theorem}[section]
\newtheorem{proposition}[theorem]{Proposition}
\newtheorem{lemma}[theorem]{Lemma}
\newtheorem{corollary}[theorem]{Corollary}
\theoremstyle{definition}
\newtheorem{definition}[theorem]{Definition}
\newtheorem{remark}[theorem]{Remark}
\newcommand{\Rea}{\operatorname{Re}}
\newcommand{\N}{\operatorname{N}}
\newcommand{\eps}{\varepsilon}
\newcommand{\sig}{\sigma_\circ}
\newcommand{\bet}{\beta_*}
\newcommand{\pos}[1]{\left(#1\right)_+}
\newcommand{\src}[1]{\nolinkurl{#1}}
\makeatletter
\g@addto@macro\UrlBreaks{\do\a\do\b\do\c\do\d\do\e\do\f\do\g\do\h\do\i\do\j\do\k\do\l\do\m\do\n\do\o\do\p\do\q\do\r\do\s\do\t\do\u\do\v\do\w\do\x\do\y\do\z\do\0\do\1\do\2\do\3\do\4\do\5\do\6\do\7\do\8\do\9}
\makeatother
\title{Optimizing compensated-probe geometry\\below seven eighths}
\author{RH--Weil research project}
\date{October 7, 2026\\[3pt]\small A derivation relative to explicitly imported results}
\begin{document}
\maketitle
\begin{abstract}
We optimize two geometric parameters in the compensated-probe argument
of the September 30, 2026 preprint \emph{The Quasi-Riemann Hypothesis}.
Assuming its specified generic arithmetic and analytic inputs and its
finite-order Hecke seven-eighths theorem, we derive the strict zero-free
half-plane
\[
 \Rea s>\frac{1507-2\sqrt{921}}{1653}
       =0.874957069799168740\ldots
\]
for finite-order Hecke $L$-functions over $\mathbb Q(\sqrt{-3})$
and, by quadratic transfer, for all Dirichlet $L$-functions.
We prove a variable-$b$ low estimate with the complete additive Gram
bound, give an exact continuous quadratic certificate for the central
high estimate, and verify the detector capacities, Euler domains and
outer rows. The critical certificate has zero additional endpoint
margin; the common gap between the family supremum and the proposed
boundary pays all real losses before the target character is chosen.
The resulting parameters are optimal for the stated low intersection,
count envelope and amplitude rectangle. A rational witness
$40773/46600$ has strictly positive additional margin.
The assertion allows the principal pole at $1$ and excludes the boundary
line. This is a cited-input derivation, not an independent certification
of the imported proof or a new Lean certificate. No improvement of a
critical-line zero proportion is asserted.
\end{abstract}
\tableofcontents

\section{Statement and exact imported hypotheses}
\label{sec:inputs}
Put $F=\mathbb Q(\sqrt{-3})$, and use the conventions of \cite{OAI}
for norm $q_a=\N(a)$, ideal and element representatives, units, ray
conditions and zero extension on nonunits. A superscript $S$ deletes
the finite Euler factors in the fixed arithmetic datum.
Our previous paper \cite{Previous} treats the boundary $69999/80000$
with $b=1/8$. We retain its imported package, but reprove every
geometric condition affected by varying $b$. Neither paper certifies
the imported package itself.

\begin{definition}[Input package $\mathcal R$]\label{def:R}
The package consists of the following results of the pinned source
\cite{OAI}, with their original coefficient classes, uniformities and
side conditions:
coefficientwise probe and Poisson identities, finite compensation,
local arithmetic tables and ray calibration; smooth separation,
Gaussian annuli, reflection sectors, generic reflected energy and the
quantitative additive Gram bound; global and buffered reciprocal
estimates, Hecke growth, functional equations and full contour tails;
actual detector witnesses, marked and plain moments, recursive
inverse moments, sixth-power amplification and uniform meshes;
fixed-ray prime asymptotics; the whole finite-order Hecke
seven-eighths zero-free theorem; and the quadratic factorization
transferring the Hecke assertion to Dirichlet functions.
\end{definition}

Define the whole-family supremum by
\[
 \bet=\sup\bigl(\{1/2\}\cup
 \{\Rea\rho:L_F(\rho,\eta)=0,\ \rho\text{ a nontrivial zero},
               \ \eta\text{ primitive of finite order}\}\bigr).
\]
The principal pole is not a zero in this definition.
The seven-eighths input is used only as $\bet\le7/8$.
In particular we do not assume the new half-plane is zero-free.
We do not directly invoke a source low or shared-interface statement
whose hypothesis fixes $b=1/8$, $\ell=1/6$ or a boundary at least $7/8$.
We instead invoke the generic results and recheck their hypotheses.
The source is frozen at commit
\src{adc7f1241b42e322a6451854ab7e4b4c146bf78a}.
All line numbers in Table~\ref{tab:inputs} refer to its September 30
\src{build/paper.tex}.

\begin{longtable}{@{}>{\raggedright\arraybackslash}p{.25\textwidth}>{\raggedright\arraybackslash}p{.25\textwidth}>{\raggedright\arraybackslash}p{.43\textwidth}@{}}
\caption{Imported inputs and the scope of their reuse.}\label{tab:inputs}\\
\toprule
Input & Source location & Conditions retained or rechecked\\
\midrule
\endfirsthead
\toprule Input & Source location & Conditions retained or rechecked\\
\midrule
\endhead
Smooth and Gaussian calculus & 1123--1209, 1288--1315 &
One joint density, weighted seminorms and whole annular tails.\\
Reflection sectors and energy & 2470--2506, 7663--7804 &
Actual residual dyads, bounded lengths, original masks and product
coefficients; frozen powerful parts counted once.\\
Additive Gram & 8340--8364 &
$Q,Y'\ge1$, fixed polynomial ranges, $P_a=Y'^2/Q\ge1$;
complete third term retained.\\
Local and full-tuple identities & 8707--9049 &
Selected primes operated on once; no division by general local
corrections.\\
Absolute bounds and tails & 9101--9195, 5672--5728 &
All heights and completed sums; fixed internal orders precede the
external tail order.\\
Detector and moment inputs & 4385--4547, 9220--9269,
12362--12389, 12492--12578 &
Actual witnesses, strict widths, finite coefficient group and
zero-capacity cases.\\
Coefficients and capacities & 15015--15446 &
Whole slots, conjugation, inducing exceptions, one shared conductor
budget.\\
Whole-bin accounting & 5808--5896, 6021--6155 &
Global move before local partitions; the same mean gain $q$.\\
Principal and outer rows & 6174--6301, 6342--6463 &
New analytic domain and $D_1(1/3)$ verified independently of the old
numerical boundary.\\
Choices and bootstrap & 16191--16463 &
Losses and mesh, then $K$, then principal $\mu$, then target and
height orders.\\
Height and continuation & 6520--6580, 400--502, 6705--6780 &
Cutoff-independent probe, common saving, supremum contradiction and
Dirichlet transfer.\\
\bottomrule
\end{longtable}

\begin{theorem}\label{thm:main}
Assume $\mathcal R$. Every finite-order Hecke $L$-function over $F$
has no zero in $\Rea s>\sig$, where
\[
 \sig=\frac{1507-2\sqrt{921}}{1653}.
\]
The same holds for every Dirichlet $L$-function, including $\zeta$.
Principal functions may have their pole at $s=1$.
The boundary line is not included.
\end{theorem}

The theorem is a relative implication under Definition~\ref{def:R}.
Its scope does not include independent verification of the source,
its dependencies or its formal proof system.
The simple critical-line proportion of \cite{AF} is not used in
the bad-row count and is not improved here.

\section{The same probe and the variable geometric scales}
\label{sec:geometry}
For fixed real parameters $\ell,b$, write
\begin{equation}\label{eq:geometry}
 l_x=\frac{1-\ell-b}{2},\quad l_y=\frac{1-\ell+b}{2},\quad
 M=1-\ell,\quad h=\frac{1+3\ell+b}{2},\quad
 X=Z^{l_x},\quad Y=Z^{l_y}.
\end{equation}
Choose a fixed even $K$, equal positive $\ell_i=\ell/K$ and
$P_i=Z^{\ell_i}$, with the original nonnegative nonzero smooth
windows in disjoint subintervals of $(1,2)$.
Keep the source $S,T,\xi$, all permitted prime representatives,
zero extensions, annular norm ratios and ray calibration.
Write $I_\eta(X,Y,Z)$ for the source completed probe and
$I_{\eta;D}$ for the same expression with $1_{D\mid cn^3}$
inserted. The marked operation in a slot is
$\overline{\eta(p)}I_{\eta;p}(X,Y,Zq_p)$; the rescaled operation
is $q_p^{-3/2}I_\eta(X/q_p,Y/q_p,Z)$, subtracted from it.
If $p_J=\prod_{i\in J}p_i$ and $p_\varnothing=1$, their
compensated product is the independent physical expression
\[
\begin{split}
 I_{\eta,\mathrm{mod}}(Z)
 =\sum_{(p_i)}\prod_iW_i(q_{p_i}/P_i)
 \sum_{J\subseteq\{1,\ldots,K\}}&
 (-1)^{|J|}q_{p_J}^{-3/2}\overline{\eta(p_{J^c})}\\
 &\cdot I_{\eta;p_{J^c}}(X/q_{p_J},Y/q_{p_J},Zq_{p_{J^c}}).
\end{split}
\]
Each slot window remains at its own $P_i$ even when another
slot rescales a completed sum.
Those completed sums are not finite; their absolute convergence is an
input in $\mathcal R$.
The compensation subsets and retained annular decompositions are finite.

The double residue at $w=1,z=1/6$ retains the exponent
\begin{equation}\label{eq:C}
 C_b(s)=s+l_x/2-1+h/6=s-\frac23-\frac b6.
\end{equation}
The Gaussian factor at this residue is still $e^{(s-5/6)^2}$.
There is no second normalizer introduced to improve the low bound.

\section{The low estimate with free b}
\label{sec:low}
\begin{proposition}\label{prop:low}
Under the generic low inputs of $\mathcal R$, for fixed
$1/6\le\ell\le1/5$ and $0<b\le3(1-\ell)/8$,
\[
 |I_{\eta,\mathrm{mod}}(Z)|
 \ll_{\eta,\mathrm{fixed},\eps}
 Z^{(1-\ell)/4-b/6+\eps}.
\]
The arbitrary exponent loss can be prescribed before the target;
constants and thresholds may depend on its fixed arithmetic data.
\end{proposition}
\begin{proof}
For a rescaled subset $J$ put
$d_J=\sum_{i\in J}\ell_i$, $\ell'=\ell-d_J$ and $M'=M-2d_J$.
Its tuple ratio $r_J=q_{p_J}/Z^{d_J}$ lies in a fixed positive
compact interval, and the common row ball is
\[
 X'=XZ^{-d_J}/r_J,\quad Y'=YZ^{-d_J}/r_J,\quad
 Q=q_{b_*}X'Y'\asymp Z^{M'},\qquad M'+\ell'-1=-3d_J.
\]
All surviving marked slots are summed inside $B_m^J$ before taking
the row norm. The common Gaussian profile and Fourier density are
fixed before separating rows and tuples.
Whole far annuli are first removed by the absolute completed-sum
input and Gaussian decay.

Here is the geometric part of the reflected-energy calculation.
Freeze the powerful part of length $O$ and the actual squarefree
dyad of length $H$. The source sectors give
\[
 \Delta_H=M'-O-H\ge-\eps_Z,\quad
 2A_0\le O+\eps_Z,\quad N_0\le A_0,\quad z_a\le\ell',
 \qquad \eps_Z=O_{\mathrm{fixed}}(1/\log Z).
\]
For retained Gaussian displacement $\theta_N$ the dual length is
$T_d=2H+2A_0+2z_a-1-\ell'-\theta_N-N_0-3B_0$.
Subtracting $H-3d_J$ gives
\[
 T_d-(H-3d_J)
 =H-M'+2A_0+2(z_a-\ell')-N_0-3B_0-\theta_N
 \le2\eps_Z+|\theta_N|.
\]
For completeness the imported reflection expression, after
counting frozen parts once, is
\[
\begin{split}
 s_{\mathrm{hyb}}&=\min\{v,z_a,(v+z_a)/3\},\qquad
 y_{\mathrm{ref}}=v+3\ell_b+e_\lambda\le T_d+\tau_{\mathrm{ref}},\\
 E_{\mathrm{ref}}&=O/2+\max(H,v+\ell_b)-S_0-B_0+z_a\\
 &\quad-s_{\mathrm{hyb}}-\ell_b-2e_\lambda/3
                  -(T_d-y_{\mathrm{ref}})_+/2.
\end{split}
\]
Here $v,\ell_b$ are nonnegative dual lengths, $S_0,B_0,N_0,A_0$
are the frozen sector lengths, and a negative $e_\lambda$ is
only $O(1/\log Z)$. Retained support gives
$\max(H,v+\ell_b)=H+O(\tau_{\mathrm{ref}}+\kappa_G+\eps_Z)$.
In this formula, the branch
$s_{\mathrm{hyb}}=z_a$ gives exponent at most $M'-\Delta_H$.
On the other branch $s_{\mathrm{hyb}}\ge v/2$, so
\[
 s_{\mathrm{hyb}}+\ell_b+2e_\lambda/3
                    +(T_d-y_{\mathrm{ref}})_+/2
 \ge T_d/4-O(\eps_Z).
\]
The exact frozen-part saving identity is
\[
\begin{split}
 O/2+\Delta_H+S_0+B_0+T_d/4
 ={}&\frac{2M'+2z_a-1-\ell'+2\Delta_H-\theta_N}{4}\\
 &+\frac{2A_0-N_0+B_0+4S_0}{4}.
\end{split}
\]
The second line is nonnegative. Using $z_a\le\ell'$ and retaining
the kernel saving before squaring and positive enlargement gives
\[
 M'+\frac{1+3\ell'-2M'-2\Delta_H+\theta_N}{4}
 +\text{small losses}.
\]
Since $1+3\ell'-2M'=5\ell-1+d_J$, both branches yield
\begin{equation}\label{eq:row}
 \sum_{\N m\ll Q}|B_m^J|^2
 \ll Z^{M'+(5\ell-1+d_J)_+/4+\eps}.
\end{equation}
This derivation does not contain $b$; the original shared primes,
empty surviving slots and powerful-part counting are unchanged.

The enlarged range still satisfies
$M'\ge1-3\ell>0$ and $l_x-d_J,l_y-d_J>0$.
Thus $Q,Y'\ge1$ eventually, in fixed polynomial ranges, and
$P_a=Y'^2/Q=q_{b_*}^{-1}Z^b\ge1$ eventually.
The generic Gram bound, including its third term, is
\[
 \sum_{\N m\ll Q}|A_m(Y')|^2
 \ll (Q/Y')(1+P_a^{1/6}+P_a^2/Y')Z^\eps
\]
with the original ray/Gauss coefficients and fixed height
seminorm cost.
Cauchy, including the separation factor $Q^{-1/2}$, bounds a
fixed tuple by
\[
 (X')^{1/2}
 Z^{\frac12\max\{0,b/6,2b-l_y+d_J\}
        +(5\ell-1+d_J)_+/8+\eps}.
\]
Tuple number, the original $q_{p_J}^{-3/2}$ coefficient and the
shortening of $(X')^{1/2}$ contribute respectively $d_J$,
$-3d_J/2$ and $-d_J/2$.
Relative to $l_x/2+b/12$, their combined additional exponent is
\[
 F(d)=-d+\frac12(11b/6-l_y+d)_+
                 +\frac18(5\ell-1+d)_+ .
\]
Every piece has slope at most $-3/8$, so its maximum is at $d=0$.
Because $\ell\le1/5$, the resulting full low exponent is
\begin{equation}\label{eq:lowfull}
 L_{\mathrm{low}}
 =\frac{l_x}{2}+\frac b{12}
       +\frac12(11b/6-l_y)_+
 =\max\left\{\frac{1-\ell}{4}-\frac b6,\frac b2\right\}.
\end{equation}
The stated range of $b$ selects the first branch.
Small reflection and separation losses, finitely many subsets,
summable joint Fourier seminorms and logarithmic dyad counts can
be absorbed into the prescribed $\eps$ in their original order.
\end{proof}

For our parameters the stronger all-subset condition
\[
 l_y-\ell-\frac{11b}{6}
 =\frac{1-3\ell}{2}-\frac{4b}{3}>0
\]
also holds. We have proved the wider low range, rather than
assuming this stronger condition is necessary.
When $b\le0$, the present direct use of the Gram input
$P_a\ge1$ is not justified; no such extension is used.
On the first branch of \eqref{eq:lowfull}, the residue intersection is
\begin{equation}\label{eq:intersection}
 C_b(\sigma)=\frac{1-\ell}{4}-\frac b6
 \quad\Longleftrightarrow\quad
 \sigma=\frac{11}{12}-\frac\ell4.
\end{equation}
Outside that branch one cannot keep this same intersection.

\section{Actual counts and the continuous endpoint certificate}
\label{sec:certificate}
Use the original legal $\kappa=3/4$, justified by $\bet\le7/8$,
and the algebraic envelope $\alpha=5/6$.
For nonfloor bins set $\delta=2a-1\in[1/50,3/4]$.
Let $q$ be the length-weighted mean amplitude over \emph{all} slots;
error and zero slots have gain zero. Put
\[
 x=q/\delta\in[0,1/2],\quad y=1/2-x,\quad
 D=\frac{37+34y}{18},\quad P=\frac{7+18y+8y^2}{9},
\]
\begin{equation}\label{eq:counts}
 J=(5/6-\delta)D+\delta P,\qquad
 R_*=1-\delta+\frac{(5/6-\delta)\delta P}{2J}.
\end{equation}
The actual inverse/plain crossing gives these functions with
$35/54\le J\le5/2$ and $1-\delta\le R_*\le1-2\delta/3$.
Section~\ref{sec:admission} verifies that the new geometric
scales meet its arithmetic prerequisites.

At \eqref{eq:intersection}, the central endpoint accounting is
\begin{equation}\label{eq:endpoint}
 E_\sigma(h)=-\frac14+\frac{5\ell}{4}+\frac b6
      +(\tfrac12+\ell)\delta+\ell\delta x-h(1-R_*).
\end{equation}
In particular $\partial_b E_\sigma=R_*/2-1/3$.
Clearing denominators gives
$-1296J E_\sigma=A(y)\delta^2-B(y)\delta+C(y)$, where
\begin{align}
 A(y)&=252+756\ell-576b+(648+288\ell+720b)y
       \notag\\[-2pt]
     &\quad +(288+1008\ell+864b)y^2+1152\ell y^3,\label{eq:ABC}\\
 B(y)&=624-1440\ell-1176b+(504-420\ell-456b)y
          +(-48+120\ell+432b)y^2,\notag\\
 C(y)&=555-2775\ell-370b+(510-2550\ell-340b)y.\notag
\end{align}
For $Q=4AC-B^2$ its constant coefficient is
\[
 Q(0)=-530496b^2+1887840b\ell-184032b
       -10465200\ell^2+678240\ell+170064.
\]
For fixed $\ell$, its maximum in $b$ occurs at
$b_{\mathrm{opt}}=(2185\ell-213)/1228$ and equals
$-1631700(1653\ell^2-66\ell-35)/307$.

\begin{proposition}\label{prop:critical}
Set
\[
 \ell_\circ=\frac{33+8\sqrt{921}}{1653},\qquad
 b_\circ=\frac{2185\ell_\circ-213}{1228}
         =-\frac4{29}+\frac{230\sqrt{921}}{26709}.
\]
Then $E_{\sig}(h)\le0$ on the entire rectangle
$\delta\in[1/50,3/4]$, $y\in[0,1/2]$.
Its unique equality point is
$y=0$, $\delta_\circ=(49-\sqrt{921})/48$.
\end{proposition}
\begin{proof}
Write $e=\ell_\circ$. It satisfies
$1653e^2-66e-35=0$ and
$1/6<e<167/1000$, $123/1000<b_\circ<124/1000$.
The coefficients of $A$ are
\[
 \frac{108036-82548e}{307},\
 \frac{160596+481716e}{307},\
 \frac{42408+781416e}{307},\
 1152e,
\]
all strictly positive. Reducing $Q$ by the same quadratic identity
gives
\begin{align*}
 Q_0&=0,\\
 Q_1&=17404800(5465-31728e)/169157,\\
 Q_2&=1200(47597384-214807491e)/169157,\\
 Q_3&=4800(131691060e-18009311)/169157,\\
 Q_4&=14400(1377691e-209998)/8903 .
\end{align*}
The isolating bounds above imply $Q_i>0$ for $i=1,\ldots,4$.
Thus, for all $y\ge0$,
\[
 A\delta^2-B\delta+C
 =A\left(\delta-\frac{B}{2A}\right)^2+\frac Q{4A}\ge0.
\]
Since $J>0$ on the actual rectangle, this proves the assertion.
Equality forces $y=0$ and
$\delta=B(0)/(2A(0))=(49-\sqrt{921})/48$,
which lies strictly between $1/50$ and $3/4$.
This is a continuous coefficient-positivity proof, not a grid
inference. Equation~\eqref{eq:intersection} gives \eqref{eq:C}
at the boundary stated in Theorem~\ref{thm:main}.
\end{proof}

\begin{proposition}[The precise parameter optimality]\label{prop:optimal}
No $\ell>\ell_\circ$ and any $b$ can make all endpoints
nonpositive while retaining \eqref{eq:intersection}, the count
\eqref{eq:counts} and the complete stated amplitude rectangle.
\end{proposition}
\begin{proof}
At the fixed point $y=0,\delta=\delta_\circ$, the count gives
$R_*=2/3$. Thus \eqref{eq:endpoint} is independent of $b$ and equals
\[
 -\frac5{12}+\frac{3\ell}{4}
          +(\tfrac12+\tfrac32\ell)\delta_\circ .
\]
It vanishes at $\ell_\circ$ and its slope in $\ell$ is
$3/4+3\delta_\circ/2>0$.
\end{proof}
This is optimality of a particular low/count certificate, not of
all possible probes or $L$-function theory. We do not claim that
actual bad rows attain the envelope or that a corresponding
zero exists.

\section{A rational witness and comparison with fixed b}
\label{sec:rational}
For $\ell_r=1/6+1/5825=5831/34950$ and $b_r=617/5000$,
\[
 l_{x,r}=\frac{2480617}{6990000},\quad
 l_{y,r}=\frac{3343183}{6990000},\quad
 h_r=\frac{1891861}{2330000},\quad \sigma_r=\frac{40773}{46600}.
\]
At these parameters $A,C,Q$ have strictly positive coefficients.
Moreover $A_0Q(y)-Q_0A(y)$ has nonnegative coefficients, and hence
\[
 -E_{\sigma_r}(h_r)\ge\frac{Q_0}{12960A_0}
 =\frac{10580567}{347281513800000}>10^{-8}.
\]
Here completing the square first gives $Q/(4A)$, and
$J\le5/2$ supplies the factor $12960$. This does not mistake
$A(0)$ for an upper bound on $A(y)$.
The exact coefficients are in the reproducible algebra record
\cite{FreeBAudit}.

The fixed-$b$ critical value in \cite{Note447} has
$\sigma_c=7/8-t_c/4$, where
$1/5842<t_c<1/5841$ is the positive root of
\[
 28224-164747520t-669772800t^2=0.
\]
Since $1/5825>1/5841$, $\sigma_r<\sigma_c$.
Also
$1653\ell_r^2-66\ell_r-35=-9289/407167500<0$;
the quadratic is increasing on $[1/6,\infty)$, so
$\ell_r<\ell_\circ$ and $\sig<\sigma_r$.
Both comparisons are exact. For orientation only,
\[
 \sig=0.874957069799168740\ldots,\qquad
 \sigma_c=0.874957200615373761\ldots .
\]
We prove the theorem at the critical algebraic value, not only
at the rational witness.

\section{Actual arithmetic admission and all row ranges}
\label{sec:admission}
For a physical row length $d$ the same-source accounting reads
\begin{equation}\label{eq:generalE}
 \begin{split}
 E_\sigma(d)={}&K_\sigma+(\tfrac12+\ell)\delta+\ell q
       -h(1-R)\\
       &+(d-h)(R+\delta/2-17/50),\qquad
 K_\sigma=2/3+\ell+b/6-\sigma .
 \end{split}
\end{equation}
Expanding the original unfactored exponent
$a-\sigma+h(17/50-1/6)-a l_y-\ell/2+\ell q
+d(R+\delta/2-17/50)$ gives this formula.
The joint bound for strict local errors and the numerator uses
one conductor deficit once; its gain remains $\ell q$.
There is no second charge for each ramified error label.

For nonfloor selected rows $1/2\le d\le h$, the slope is at least
$33/50-\delta/2\ge57/200>0$, so the endpoint controls the whole range.
The floor $a=51/100$ uses $\delta=1/50$, $R=1$, $q\le1/100$
independently of a witness.
The no-slot count on $1/100\le d\le1/2$ uses
$R=76/75-2\delta/3$, covering the floor, with positive slope
$101/150-\delta/6$.
At $\ell=e=\ell_\circ,b=b_\circ$, direct calculation gives
\begin{align}
 l_y-e-11b/6&=(1347-7133e)/1842>2/25,\notag\\
 5e-h&=(6411e-1015)/2456>1/50,\label{eq:margins}\\
 E_{\mathrm{floor}}(h)&=(290401e-49533)/184200<-1/200,\notag\\
 E_{\mathrm{middle}}&=(292151e-84703)/1473600<-1/50.\notag
\end{align}
The middle bound follows from $q\le\delta/2$:
its value at $d=1/2$ is
$-73/300+(13/50)\ell-(49/300)b
+\delta(1/6+3\ell/4-b/4)$,
maximized at $\delta=3/4$.

Assume for contradiction $\bet>\sig$, and put
$\Delta=\bet-\sig>0$. The bootstrap gives
$\Delta\le(e-1/6)/4$. Choose $\zeta=\Delta/32$.
Then $\zeta<1/384000$ and \eqref{eq:margins} implies
\[
 \ell/(h+\zeta)>1/5>7/37,\qquad h+\zeta<1.
\]
The selected-row slope is at most $2$, so the extension from $h$
to $h+\zeta$ costs at most $\Delta/16$.
The actual comparison is with $C_b(\bet)$:
\begin{equation}\label{eq:Delta}
 E_{\bet}(d)=E_{\sig}(d)-\Delta.
\end{equation}
Proposition~\ref{prop:critical} thus provides central saving at
least $15\Delta/16$ before small real losses.
Zero additional margin at the critical endpoint does not preclude
continuation.

We spell out the moment prerequisites that varying $b$ must preserve.
The adaptive parameter is $t=1+\delta P/(2J)$.
The original witnesses have $r+m\ge t-o(1)$, $r\le t+o(1)$,
$m\le1/2+o(1)$ and the original squared spike sizes.
Writing $x=q/\delta$, the original inverse/plain threshold is
\[
 r_*(t)=\frac{(2-8x/9)t-5x/9}{3-17x/9},
 \qquad r_*(t)\ge23/37,\quad 1\le t\le3/2 .
\]
On the inverse branch $r\ge r_*(t)$, $r<1$ and positive
$z_M=(1-r)/2$. For a fixed capacity decrement $\nu_0>0$, requesting
$z\le z_M-\nu_0$ gives the two inverse widths
\[
 1-r-2z\ge2\nu_0,\qquad
 3-2r-8z\ge8\nu_0+9/37-o(1).
\]
On the plain branch $r\le r_*(t)$, $m<1/2$ and positive
$z_P=(1-2m)/(6\kappa)$; requesting $z\le z_P-\nu_0$ gives,
with $m\ge1/3-o(1)$,
$1-2m-6\kappa z\ge6\kappa\nu_0$.
Maximum demands are $z_M\le7/37$ and $z_P\le2/27+o(1)$.
The supply above therefore suffices; it does not itself strengthen
the moment theorem.
Long-witness and zero-capacity endpoints retain the original
zero-slot or sixth-power cases.

Only $d\ge1/2$ uses selected primes, so their physical lengths
$w_i=\ell_i/d\le2\ell_i$.
Choose losses, decrement, mesh and rounding before $K$;
then every $w_i$ can be smaller than the fixed plain mesh.
The greedy selection loses at most one whole-slot length and
uses the gain averaged over all slots.
Coefficients remain the fixed finite $\Theta$ combination of the
original arithmetic phases; all witnesses and selected factors
are conjugated together on the plain side.
The physical slots remain disjoint from the internal plain pool.
Outside-$S$ ramification excludes the original inducing exceptions;
the remaining $S$-supported sixth-power-free rows are finite and
belong to the independently handled bounded rows.
These are precisely the prerequisites of the unchanged count
\eqref{eq:counts}.

\section{Euler domains, principal residue and outer rows}
\label{sec:contours}
Varying $b$ does not alter a local coefficient or zero extension.
The full selected-prime operation and the coefficientwise
physical identity remain the source identities.
The resulting raw Mellin weight is
\[
 X^{1/2-z}Z^{s+z-1}Y^{w-1}e^{(s+z-1)^2}
 M(z)\widehat W_1(w).
\]
The full tuple is the finite selected-prime sum with selected
factors $G_p$ and the normally convergent unselected correction
product. It is never defined by dividing by a general $H_p$.

The same local finite-table calculation as in \cite{Previous}
applies on the enlarged principal box:
\[
 \Rea s\ge\sig,\qquad
 \Rea w\ge19/20,\qquad \Rea z\ge33/200.
\]
Indeed $\sig>5/6>401/600$. The good and ramified defects are
\begin{align*}
 c_{\mathrm{good}}
 &=\min\{6\sig-401/100,\sig-3/50,47/50\}
   =\sig-3/50>81/100,\\
 c_{\mathrm{ram}}
 &=\min\{\sig-1/20,3\sig-3/2\}
   =\sig-1/20>82/100.
\end{align*}
The denominators $1-R,1-V,1-D$ stay away from zero.
The general good local term with exponent $4-6s-6z$ retains
the correction $\theta=(-\Rea w)_+$.
In $D_1(\eps_0)$ the exponent used for the correction is
$\eps_H=\min(\eps_0,1/50)$, not always $\eps_0$.
These strict margins imply local normal convergence and full
height majorants, not just values on a single contour.

Whole-bin contours are moved on the global line before amplitude
and error partitions, followed by the original buffered local move.
All strict labels use the joint numerator estimate once.
The principal correction is the original normally convergent product
\[
 H_\eta(s)=\prod_{p\notin S}H_p(s,1,1/6)
          =\mathcal H_{\eta,1}(s,1,1/6),
\]
with the same final $S$ as the physical expression.
For the principal row choose a target-independent $P_0$ such that
the positive product majorant gives $|H_\eta(s)-1|\le1/2$ on
$\Rea s\ge\sig$. Enlarging the final $S$ only shrinks this majorant.
More explicitly, put
\[
 S_i(Z)=\sum_{p\in\mathcal P_i(Z)}
 W_i(q_p/P_i)q_p^{-5/6}
 \sim c_iP_i^{1/6}/\log P_i,\qquad
 A_T=(-1)^K Z^{-\ell/6}\prod_i S_i(Z).
\]
The fixed-ray calibration supplies $c_i>0$, the same
$c_S=\widehat W_1(1)M(1/6)
(\operatorname{Res}_{v=1}\zeta_F^S(v))^2/6>0$,
and eventually nonzero $|A_T|\asymp(\log Z)^{-K}$,
with subpower inverse. These are the same finite $S$ and
actual ray slots used in the physical expression.

The principal contours start with global $s$, then pass the
residues at $w=1,z=1/6$, with final nonresidue lines $19/20$ and
$33/200$. The double-residue $s$ line is moved only to the right,
to $2$. The three error margins are
\[
 m_w=l_y/20>0,\qquad m_z=h/600>0,\qquad
 0<\mu<\sig\min_i\ell_i.
\]
The four principal local error exponents remain
$-\Rea s,-6\Rea z,4-5\Rea s-6\Rea z,1-\Rea w-6\Rea z$;
their weakest decay at the residue is $\sig$.

Small rows $d<d_{\min}=1/100$ and bounded nontrivial units are
handled separately. The old numerical invocation of $D_1(3/8)$
is replaced by the local proof on $D_1(1/3)$, valid since
$\sig+1/2>4/3$.
The full absolute tuple and numerator bound give, relative
already to $C_b(\bet)$,
\[
 E_{\mathrm{small}}
 =h(13/75)-l_y/2+63/5000
 =\frac{2020475e-1127329}{9210000}<-2/25.
\]
No additional $\Delta$ is charged to this formula.
Large rows $U>Z^{h+\zeta}$ use the original full absolute sum
on $(2,2,z_\infty)$:
\[
 Z^{B_0+(h+\zeta)(1+\eps)-\zeta z_\infty+\eps},
 \qquad B_0=l_x/2+1+l_y=7/4-3\ell/4+b/4.
\]
After fixing $\zeta$, choose a fixed sufficiently large
$z_\infty$ before the target; then small losses and all internal
height orders. This gives any prescribed saving without changing
an order with a moving row or with $Z$.

\section{Uniform choices and whole-family continuation}
\label{sec:continuation}
Fix the critical geometry and the common contradiction gap
$\Delta$ before any target character.
Use small fixed fractions of $\Delta$ for moment losses, strict
capacity decrement, uniform mesh and rounding.
Then choose even $K$ and the disjoint windows.
Only afterwards put $\mu=\sig\ell/(2K)>0$ and choose amplitude
widths, prime losses and the small contour displacement.
Principal losses are separately smaller than
$m_w/4,m_z/4,\mu/4$; central total real losses are below
$\Delta/4$, and floor, middle and small losses leave half their
respective independent margins.
There is no circular selection of $K$ using a $\mu$ that
depends on $K$.
Fix $\zeta$ and $z_\infty$ in the same target-independent phase.

Let
\[
 m_0=\min\{\Delta,m_w,m_z,\mu,1/200,1/50,2/25\}>0,\qquad m=m_0/4.
\]
The common low loss, including the inverse normalizer, is
$\omega=\Delta/2<\Delta$.
For a chosen target, fix its arithmetic datum and final $S$,
and use that same datum for every physical expression,
calibration and correction.
For sufficiently large $Z$ define the cutoff-independent
normalized physical function $J_\eta$; for all $Z>0$ define $f_\eta$:
\[
 J_\eta(Z)=I_{\eta,\mathrm{mod}}(Z)/(c_S A_T(Z)),\qquad
 f_\eta(Z)=\frac1{2\pi i}\int_{(2)}
  Z^{C_b(s)}e^{(s-5/6)^2}\frac{H_\eta(s)}{L_F^S(s,\eta)}\,ds.
\]
The low bound gives
$|J_\eta|\ll_\eta Z^{C_b(\sig)+\omega}$.
The principal and all high estimates, with the same $S$ and
normalizer, give finite $A_\eta,B_\eta$ such that
\begin{equation}\label{eq:contract}
 |J_\eta-f_\eta|\ll_{\eta,N}
 Z^{C_b(\bet)-m}(1+T_1)^{A_\eta}
       +Z^{B_\eta}T_1^{-N}.
\end{equation}
Neither function contains $T_1$.
The internal moment, Sobolev and seminorm orders giving
$A_\eta,B_\eta$ are fixed before the external $N$.
All retained frequencies share one cumulative $T_1/2$ allocation.

For the actual detector ceiling $\tau_{0,\eta}>0$ choose
\[
 0<\tau_\eta\le
 \min\{d_{\min}/100,\tau_{0,\eta},m/[4(A_\eta+1)]\},
 \qquad T_1=Z^{\tau_\eta}.
\]
After this choose $N_\eta$ so that
$B_\eta-N_\eta\tau_\eta<C_b(\bet)-m/2$, and only then the
threshold in $Z$.
Equation~\eqref{eq:contract} has a common target-independent
saving $\sigma_{\mathrm{hi}}=m_0/8>0$.
The target dependence of $\tau_\eta$, $N_\eta$, constants and
thresholds does not change this saving.

\begin{proof}[Proof of Theorem~\ref{thm:main}]
All prerequisites for the source continuation argument have now
been re-established at the new geometry.
For clarity we describe its contradiction step.
The Gaussian and rightward shifts give rapid decay of $f_\eta$
as $Z\to0$. No small-$Z$ definition of $J_\eta$ or inverse
$A_T$ is needed.
The common high saving and the strict low gap imply locally
uniform convergence of the Mellin/Fourier identification to a
line left of $\bet$, by the common amount
$\eps_*=\min\{\Delta/2,m_0/8\}>0$.
Specifically,
$\mathcal F_\eta(s)=\int_0^\infty
f_\eta(Z)Z^{-C_b(s)}\,dZ/Z$ converges locally uniformly on
$\Rea s>\bet-\eps_*$. On the initial line $2$ the Gaussian
Fourier identification gives
$\mathcal F_\eta(s)=e^{(s-5/6)^2}H_\eta(s)/L_F^S(s,\eta)$.
Analytic continuation, without moving a contour across a target
zero, therefore makes
$H_\eta/L_F^S(\,\cdot\,,\eta)$ holomorphic to that line.
The positive product choice $|H_\eta-1|\le1/2$ excludes
cancellation by the numerator.

The supremum $\bet$ need not be attained: its definition gives
a target zero with real part greater than $\bet-\eps_*$.
All the preceding choices are common before that target, and
the resulting holomorphic reciprocal contradicts this zero.
Thus $\bet\le\sig$.
Finite Euler deletions have no zero for $\Rea s>0$.
The quadratic Dirichlet factorization in $\mathcal R$ transfers
the assertion to all Dirichlet functions.
At $s=1$ retain the principal pole and use
$L(1,\chi_{-3})=\pi/(3\sqrt3)\ne0$ for the zeta factor.
\end{proof}

\section{Relation to zero proportions and further work}
The argument does not use the simple critical-line proportion
of Alp\"oge--Furman \cite{AF}. Its count improvement is geometric
within a different actual family and moment calculation.
Conversely, a narrower global zero-free strip alone does not pay
the signed four-point correlation in the original finite
critical-line response.
Our companion note \cite{Note448} keeps the natural masks,
ghost poles and finite frame while proving a canonical
Type-I admission and low-prime fourth-norm payment.
Its conditional one-factor contour estimate is weaker than a
local unconditional finite-ratio Gram bound; neither supplies a
new fourth-moment constant or a larger simple-zero proportion.

Proposition~\ref{prop:optimal} identifies the next boundary task:
one must strengthen the actual count/gain relation or the
admissible moment mechanism near $y=0,\delta=\delta_\circ$,
rather than only changing $b$.
For proportions the remaining task is the original joint
signed response, including the common centering, cross terms,
all relevant cells and finite-frequency alias bookkeeping.
Neither unfinished task is replaced by the optimality statement.

\appendix
\section{Exact algebra and verification scope}
\label{app:verification}
The pinned source canonical LF SHA-256 is
\[
 \text{\src{42a5ee0febca59fd1def55cfd6c6808c322ef4237d7726303ea52f711deac6a3}}.
\]
The derivation is recorded in note 449 \cite{Note449}.
The independent parameter report and exact script \cite{FreeBAudit}
use rational arithmetic and $\mathbb Q(\sqrt{921})$, with no
floating-point inference. The script checks the displayed
discriminant coefficients, positive signs, equality point,
rational witness and geometric margins. Its 841 checks include
189 direct endpoint cases as implementation checks.
The continuous proof is the polynomial identity and positivity
argument in Section~\ref{sec:certificate}, not those cases.

Run, from the repository root,
\begin{quote}
\texttt{python}\ \src{scripts/hybrid_free_b_geometry_exact_audit.py}
\end{quote}
The script prints its own canonical hash and an exact JSON record.
Final source, PDF and independent full-paper review hashes are
bound separately in the build manifest. The script does not
certify actual prime sums, imported estimates, infinite contours,
the source proof kernel or RH.
The earlier $69999/80000$ paper is preserved as its delivered
version; the present result is a separate paper.

\begin{thebibliography}{9}
\bibitem{OAI}
OpenAI, \emph{The Quasi-Riemann Hypothesis},
September 30, 2026 preprint, pinned source at
\href{https://github.com/openai/math/blob/adc7f1241b42e322a6451854ab7e4b4c146bf78a/preprints/The-Quasi-Riemann-Hypothesis-September-30-2026/build/paper.tex}
{OpenAI/math, commit adc7f1241}.
\bibitem{Previous}
RH--Weil research project,
\emph{A small improvement of the seven-eighths zero-free boundary
from a compensated probe}, October 7, 2026.
\src{papers/seven-eighths-boundary-improvement-paper.tex}.
\bibitem{Note447}
RH--Weil research project, Note 447, October 7, 2026.
\src{notes/447-larger-slot-discriminant-and-critical-strip.md}.
\bibitem{Note448}
RH--Weil research project, Note 448, October 7, 2026.
\src{notes/448-canonical-type-i-admission-finite-gram-and-low-prime-fourth-norm.md}.
\bibitem{Note449}
RH--Weil research project, Note 449, October 7, 2026.
\src{notes/449-free-b-compensated-geometry-and-optimal-relative-boundary.md}.
\bibitem{FreeBAudit}
RH--Weil research project, Free-$b$ optimization and exact audit,
October 7, 2026.
\src{reviews/2026-10-07/hybrid-free-b-geometry-optimization.md};
\src{scripts/hybrid_free_b_geometry_exact_audit.py}.
\bibitem{AF}
L. Alp\"oge and R. Furman,
\emph{More than two thirds of the zeros of the Riemann zeta function
are simple and on the critical line}, arXiv:2608.13637v2, 2026,
\href{https://arxiv.org/html/2608.13637v2}{version 2}.
\end{thebibliography}
\end{document}
